\begin{titlepage}
\thispagestyle{empty}
\centering
{\large\bfseries UNIVERSIDAD NACIONAL DEL ALTIPLANO\par}
\vspace{3mm}
{\small FACULTAD DE INGENIERÍA MECÁNICA ELÉCTRICA, ELECTRÓNICA Y SISTEMAS\par}
{\small ESCUELA PROFESIONAL DE INGENIERÍA MECÁNICA ELÉCTRICA\par}
\vspace{8mm}
\includegraphics[height=27mm]{figuras/escudo-unap.png}\par
\vspace{8mm}
{\large\bfseries CENTRALES HIDROELÉCTRICAS\par}
\textcolor{azul}{\rule{\textwidth}{1pt}}\par
\vspace{5mm}
{\LARGE\bfseries\color{azul} MINICENTRAL HIDROELÉCTRICA LIMBANI\par}
\vspace{4mm}
{\large Objetivo de estudio en punto de conexión: 150 kW\par}
\vspace{5mm}
{\large\bfseries ANTEPROYECTO DE PRIMERA UNIDAD\par}
{\normalsize Memoria, base normativa, cálculos justificativos,\\especificaciones y esquemas conceptuales\par}
\vspace{9mm}
\begin{tabular}{L{4.1cm}L{9cm}}
\textbf{Alumno} & Aquiles Taylor Ramos Yapo\\
\textbf{Código de estudiante} & 227368\\
\textbf{Responsable de asignatura} & MEZA MAROCHO GREGORIO\\
\textbf{Sección} & GRUPO A\\
\textbf{Localización propuesta} & Limbani, Sandia, Puno\\
\textbf{Revisión} & R02 - 13 de octubre de 2026\\
\textbf{Nivel del estudio} & Anteproyecto académico con hipótesis explícitas\\
\end{tabular}
\vfill
\textcolor{azul}{\rule{\textwidth}{.5pt}}\par
{\small PUNO - PERÚ | 2026\par}
\end{titlepage}
